\documentclass[12pt]{article}
\usepackage[T1]{fontenc}
\usepackage[utf8]{inputenc}
\usepackage{lmodern}
\usepackage{textcomp}
\usepackage{enumitem}
\usepackage{geometry}
\usepackage{graphicx}
\usepackage{amsmath, amssymb}
\geometry{margin=1in}
\begin{document}
\begin{center}
Health Sciences - Graduate Level Assessment 21 \
Total Marks: 40 \
Answer all questions.
\end{center}

\section*{Multiple Choice (10 marks)}
\begin{enumerate}[label=\arabic*.]
\item The palatine processes begin to fuse at
\begin{enumerate}[label=(\alph*)]
    \item eight weeks post-fertilization.
    \item fourteen weeks post-fertilization.
    \item seven weeks post-fertilization.
    \item twelve weeks post-fertilization.
    \item four weeks post-fertilization.
\end{enumerate}
\item The main reason given by men and women for their long lasting marriage is that
\begin{enumerate}[label=(\alph*)]
    \item They have an arranged marriage
    \item They agree on almost everything
    \item They have a strong religious belief
    \item They maintain separate social lives
    \item They always take vacations together
\end{enumerate}
\item What is a tracheostomy?
\begin{enumerate}[label=(\alph*)]
    \item An opening in the anterior chest wall.
    \item An opening in the anterior wall of the trachea below the cricoid cartilage.
    \item An opening in the anterior chest wall, above the diaphragm.
    \item An opening in the posterior wall of the esophagus.
    \item An opening in the posterior wall of the trachea above the cricoid cartilage.
\end{enumerate}
\item The coronary arteries
\begin{enumerate}[label=(\alph*)]
    \item arise from the ascending aorta and do not fill during either systole or diastole.
    \item arise from the arch of the aorta and do not fill during either systole or diastole.
    \item arise from the ascending aorta and fill during systole.
    \item arise from the arch of the aorta and fill during systole.
    \item arise from the pulmonary artery and fill during systole.
\end{enumerate}
\item What is the primary lipoprotein secreted from the liver that is at least partially composed of dietary derived lipids?
\begin{enumerate}[label=(\alph*)]
    \item Lp(a)
    \item HDL
    \item ApoB-100
    \item IDL
    \item VLDL
\end{enumerate}
\end{enumerate}

\section*{True / False (10 marks)}
\textit{Answer True or False}
\begin{enumerate}[label=\arabic*.]
\setcounter{enumi}{5}
\item True or False: The patellar tap is an effective clinical test for assessing knee effusion.
\item True or False: The best available evidence for an association between sugars and dental caries risk comes from case-control studies.
\item True or False: An enlarged submandibular salivary gland is typically palpable both intraorally and extraorally during a clinical examination.
\item True or False: Checking the pulse on both sides of the neck is not necessary when examining the jugular venous pulse.
\item True or False: Aspartic acid has been replaced by phenylalanine in the molecular report denoted as Asp 235 Phe.
\end{enumerate}

\section*{Long Form Response (20 marks)}
\textit{Answer each question with a clear and complete explanation.}
\begin{enumerate}[label=\arabic*.]
\setcounter{enumi}{10}
\item Discuss the emergence of MERS as a significant coronavirus, analyzing its impact on global health, the factors contributing to its spread, and the challenges in managing such emergent viruses.
\item Discuss the pattern of inheritance of glucose-6-phosphate dehydrogenase (G 6 PD) deficiency, analyzing its implications for genetic counseling and population health.
\end{enumerate}

\vfill
\noindent\textit{End of Paper}
\end{document}